\begin{tikzpicture}[font=\small,node distance=8mm and 9mm]
\node[panel,text width=4cm](a){\textbf{Datos fuente}\\$B_j$ y $f_j$\\en cada $E_j$};
\node[proceso,text width=4cm,right=of a](b){\textbf{Producto nodal}\\$C_j=B_j f_j$\\antes de interpolar};
\node[proceso,text width=4cm,right=of b](c){\textbf{Promedio de Eq. 13}\\Interpolar $C(E,s)$\\$C(Q,s)=\langle C\rangle_w$};
\draw[flujo](a)--(b);\draw[flujo](b)--(c);
\node[aviso,text width=13.4cm,below=12mm of b](ex){\textbf{Ejemplo mínimo con dos energías y pesos iguales}\\$(B_1,f_1)=(1,2)$, $(B_2,f_2)=(3,4)$:\quad $\langle Bf\rangle=(2+12)/2=7$.\\Multiplicar promedios da $\langle B\rangle\langle f\rangle=2\times3=6$.\\La diferencia $7-6=1$ es la covarianza ponderada, no una incertidumbre de medida.};
\node[nota,text width=13.4cm,below=6mm of ex]{Valores algebraicos didácticos, sin significado dosimétrico. Las ecuaciones 11, 12 y 13 no usan necesariamente los mismos pesos.};
\end{tikzpicture}
